\documentclass[10pt,a4paper]{amsart}
\usepackage[margin=19mm]{geometry}
\usepackage{amsmath,amssymb,mathtools}
\usepackage[scr=rsfs,cal=euler]{mathalfa}
\usepackage{xfrac}
\usepackage[unicode,hidelinks,bookmarks=false]{hyperref}
\newcommand{\ie}{i.e.\ }
\newcommand{\cf}{cf.\ }
\newcommand{\resp}{respectively\ }
\newcommand{\Subsection}{\subsection}
\providecommand{\nobreakdash}{\nobreak\hbox{-}}
\setlength{\emergencystretch}{2em}
\multlinegap0pt
\usepackage{float}
\usepackage{enumerate}
\usepackage[normalem]{ulem}
\usepackage{amssymb}
\newtheorem{lemma}{Lemma}
\usepackage{cleveref}
\crefname{lemma}{Lemma}{Lemmas}
\title{On the equivariant $KU_G$-local sphere for finite abelian groups}
\author{Yingxin Li}
\date{Source: arXiv:2605.30285v3 (CC BY 4.0)}
\begin{document}
\maketitle

\noindent\emph{Source and licence.} On the equivariant $KU_G$-local sphere for finite abelian groups. Version arXiv:2605.30285v3 (CC BY 4.0), distributed by arXiv under the Creative Commons Attribution 4.0 International licence, http://creativecommons.org/licenses/by/4.0/, as recorded in the arXiv licence metadata for this exact version and shown on the abstract page https://arxiv.org/abs/2605.30285v3. Full licence notice: \texttt{licenses/arxiv-2605.30285-CC-BY-4.0.txt}.\par\smallskip

\noindent\textit{Editorial scope.} Counted unit: the article's lemma inflation (source0.tex lines 622--628). The article's complete original proof of it is delivered with it (source0.tex lines 629--685, one \texttt{proof} environment of 57 lines). No mathematics is written, repaired or re-derived: the statement and the proof are the article's own lines, in the article's own notation, and no retained line is substituted or re-flowed.  Retained uncounted as the source-local context the proof needs: lines 589--591.  Nothing was omitted from any retained span, so the package carries no omission record.  No retained line contains a citation, so the wrapper carries no bibliography and no bibliography entry.  No retained line contains a figure environment, an image-inclusion command or drawing code, so no image is delivered and the package declares no asset dependency.  Editorial formatting only, in addition: an AMS \texttt{amsart} preamble that restates the article's own declarations and macros used by the retained lines, namely the environments \texttt{lemma} on separate counters, together with the standard packages those lines need.  The retained environments print in this document as Lemma 1 in order of appearance, whereas the article prints the same items with the numbering of its own full document.  The complete native source of the article as distributed in the arXiv e-print for this exact version is bundled unchanged as \texttt{sources/201664/source0.tex}.
\begin{lemma}\label{lem:fixed point of suspension spectra}
    Let $H$ and $K$ be finite groups with coprime orders, let $G=H\times K$, and let $X\in \mathcal{S}^K_*$ be a pointed $K$-space. Regard $X$ as a $G$-space via the quotient map $G\to K$. Then, as an $H$-spectrum, $(\Sigma^\infty_G X)^K\simeq \operatorname{Inf}_e^H (\Sigma_K^\infty X)^K$. 
\end{lemma}

\begin{lemma}\label{lem:fixed point of ECyc and inflation}
    For any subgroup $H\subset G$, let $P=H\cap N_p$ and $L=H\cap N$, then for any $N_p$-spectrum $E$,
    \[ 
    (ECyc_+\wedge \operatorname{Inf}_{N_p}^G E)^H\simeq (ECyc^P_+ \wedge \operatorname{Res}_P^{N_p} E)^P \wedge (\bigvee_{(T)\in Cyc^L/L}BW_L T_+).
    \]
    Here $Cyc^P$ (resp., $Cyc^L$) on the right-hand side of the equivalence is the family of all cyclic subgroups of $P$ (resp., $L$), $(T)$ is the conjugacy class of $T$.
\end{lemma}

\begin{proof}
   For any $K\subset G$, let $Cyc^K$ be the family of all cyclic subgroups of $K$, then we have 
   \[
   \operatorname{Res}_K^G ECyc_+\simeq ECyc_+^K.
   \] 
   Since $p\nmid |N|$, for any cyclic subgroup $P_0\subset N_p$ and $L_0\subset N$, $P_0\times L_0$ is also a cyclic subgroup. By the definition of $ECyc$, we have
   \[
   ECyc_+\simeq ECyc^{N_p}_+\wedge ECyc^N_+.
   \]
    Since $H=P\times L$, there is an equivalence of spectra $X^H\simeq (X^L)^P$ for any $G$-spectrum $X$. 
    
    Consider the $L$ fixed point of $ECyc_+\wedge \operatorname{Inf}_{N_p}^G E$ as a $P$-spectrum, we have
    \[
    \begin{aligned}
        (ECyc_+\wedge \operatorname{Inf}_{N_p}^G E)^L & \simeq (ECyc^L_+\wedge \operatorname{Inf}_{P}^{H} (ECyc^P_+\wedge \operatorname{Res}_{P}^{N_p}E))^L \\
        & \simeq (ECyc^P_+\wedge \operatorname{Res}_{P}^{N_p}E) \wedge (\Sigma_H^\infty ECyc_+^L)^L.
    \end{aligned}
    \] 
    By \Cref{lem:fixed point of suspension spectra}, 
    \[
    (\Sigma_H^\infty ECyc_+^L)^L\simeq \operatorname{Inf}_e^{P}(\Sigma_L^\infty ECyc_+^L)^L,
    \] 
    thus
    \[
    \begin{aligned}
        (ECyc_+\wedge \operatorname{Inf}_{N_p}^G E)^H&\simeq ((ECyc_+\wedge \operatorname{Inf}_{N_p}^G E)^L)^P \\
        &\simeq ((ECyc^P_+\wedge \operatorname{Res}_{P}^{N_p}E) \wedge (\Sigma_H^\infty ECyc_+^L)^L)^P\\
        &\simeq ((ECyc^P_+\wedge \operatorname{Res}_{P}^{N_p}E) \wedge \operatorname{Inf}_e^{P}(\Sigma_L^\infty ECyc_+^L)^L)^P\\
        &\simeq (ECyc^P_+\wedge \operatorname{Res}_{P}^{N_p}E)^P \wedge (\Sigma_L^\infty ECyc_+^L)^L,
    \end{aligned}
    \]
    where the last equivalence follows from the formula 
    \[
    (\operatorname{Inf}_e^P x\cdot y)^P=x\cdot y^P, \quad \forall \ x\in \mathrm{Sp}, \ y\in \mathrm{Sp}^P.
    \]
    By tom-Dieck splitting formula, 
    \[
    (\Sigma_L^\infty ECyc_+^L)^L\simeq \bigvee_{(T)\subset L}\Sigma^\infty EW_L T_+\wedge_{W_L T} (ECyc^L_+)^T.
    \] 
    If $T\subset L$ is not cyclic, $(ECyc^L_+)^T$ is $W_L T$-equivariant contractible, and 
    \[ EW_L T_+\wedge_{W_L T} (ECyc^L_+)^T\simeq \ast.\] 
    If $T\subset L$ is cyclic, then for any subgroup $K \subset W_L T$,
    \[
    (EW_L T_+\wedge (ECyc_+^L)^T)^{K}\simeq\begin{cases}
        S^0 & K=\{e\},\\
        \ast & K\neq \{e\},
    \end{cases}
    \]
    which implies that there is an equivalence of $W_L T$-spaces $EW_L T_+\wedge (ECyc_+^L)^T\simeq EW_L T_+$, so 
    \[
    EW_L T_+\wedge_{W_L T} (ECyc^L_+)^T\simeq BW_L T_+.
    \] 
    Therefore, we have $(\Sigma_L^\infty ECyc_+^L)^L\simeq \bigvee_{(T)\in Cyc^L/L}BW_L T_+$, and 
    \[  
    (ECyc_+\wedge \operatorname{Inf}_{N_p}^G E)^H \simeq (\operatorname{Res}_P^{N_p} E\wedge ECyc_+)^P \wedge (\bigvee_{(T)\in Cyc^L/L}BW_L T_+).
    \]
\end{proof}

\end{document}
